\section{Expressions and evaluation}
\label{sec:evaluation}

Evaluation is strict and left-to-right; every operand is evaluated
exactly once per execution of the enclosing expression. Side effects
include calls, reads of places with side-effecting subexpressions, and
the mutating array operators.

\subsection{Numeric semantics}

\ic{long} arithmetic is two's complement modulo $2^{64}$: overflow
\emph{wraps}, it does not trap. Division and remainder truncate toward
zero (C semantics) and are subject to:

\begin{itemize}[nosep]
\item division or remainder by zero is a runtime error
  (Section~\ref{sec:errors});
\item \ic{LONG\_MIN / -1} evaluates to \ic{LONG\_MIN} (the quotient does
  not fit and would fault on x86) and \ic{LONG\_MIN \% -1} evaluates to
  \ic{0}. Implementations must special-case these two.
\end{itemize}

\ic{float} arithmetic is IEEE-754 binary64. Mixed long/float operands
widen the long side. Unary minus on a \ic{float} flips the sign, so
\ic{-0.0} is a negative zero (\ic{1.0 / -0.0} is $-\infty$) and
\ic{-(-x)} is \ic{x} for every \ic{x}.

\subsection{Logic}

\ic{\&\&} and \ic{||} short-circuit: the right operand is evaluated
only when the left does not decide the result. Operands are not
required to be \ic{bool} --- any non-zero value is true, and the result
of either operator is a \ic{bool}.

\subsection{Strings}

Strings are immutable values: concatenation \ic{a + b} produces a new
string. If exactly one side of \ic{+} is a string, the other side is
stringified (Section~\ref{sec:builtin-tostring}) and concatenated.
A format string, \ic{f"\{x\} of \{n:>4\}"}, writes such a concatenation
with fields (Sections~\ref{sec:literals} and~\ref{sec:format-spec}).
Comparisons \ic{<} \ic{<=} \ic{>} \ic{>=} on two strings compare bytes
lexicographically; \ic{==} compares contents.

\paragraph{Appending.}\label{sec:str-append} \ic{s += t} means \ic{s = s + t}, but a loop of
appends is linear, not quadratic: when \ic{s} is a variable whose
address is never taken and that no closure captures, and \ic{s} still
holds the string its own previous \ic{+=} built, the bytes of \ic{t} are
written in place behind it (the buffer doubles when full). Any other
use of \ic{s} --- a read, a copy into another binding, an element or
argument --- ends that ownership, and the next append copies once, so
an append never changes a string that anything else can see. Building
a string with \ic{+=} costs amortized $O(1)$ per appended byte, like
pushing onto an array.

\subsection{Assignment}
\label{sec:assignment}

\ic{p = e} evaluates \ic{e} and stores it into the place \ic{p}.
Assignment is a statement form, not an operator: it cannot appear inside
another expression, so \ic{let y = (x = 7);} and \ic{let y = x = 7;} are
syntax errors (so are the compound forms of
Section~\ref{sec:exactly-once}; only \ic{++} and \ic{--} nest). As the
last statement of a block it gives the block its value, the value
stored: \ic{let y = \{ x = 5 \};} sets both \ic{x} and \ic{y} to 5, and
a function body ending in an assignment returns the assigned value.
Section~\ref{sec:references}
defines what storing into arrays, records and variables means for
aliasing.

\subsection{Increment, decrement and compound assignment}
\label{sec:exactly-once}

The four increment/decrement forms and the five compound-assignment
operators share one rule:

\begin{quote}
\textbf{Exactly-once rule.} The target place's subexpressions are
evaluated \emph{exactly once}. Writing \ic{a[f()] += e} evaluates
\ic{f()} one time, no matter how the place is read and written.
\end{quote}

\noindent Formally, let a \ty{place} be $p ::= x \mid a[i] \mid o.f$,
possibly nested. Define $\mathrm{cap}(p)$, the capture of a place into
fresh variables (every non-name subexpression runs once):

\begin{align*}
\mathrm{cap}(x) &= \epsilon \\
\mathrm{cap}(a[i]) &= \mathrm{cap}(a);
   \;\; \text{if } a \text{ is not a name: } \texttt{let } \beta = a \texttt{;}
   \;\; \texttt{let } \iota = i \texttt{;} \\
\mathrm{cap}(o.f) &= \mathrm{cap}(o);
   \;\; \text{if } o \text{ is not a name: } \texttt{let } \beta = o \texttt{;}
\end{align*}

\noindent with $\mathrm{read}(p)$ and $\mathrm{write}(p, v)$ the
corresponding read and store through the captured names. Then the
\emph{numeric} forms desugar as:

\begin{align*}
\texttt{p++} &\;\equiv\; \mathrm{cap}(p);\;\texttt{\{}~\texttt{let t = read}(p)\texttt{;}
   ~~\texttt{write}(p, \texttt{t + 1})\texttt{;}~~\texttt{t}~\texttt{\}} \\
\texttt{++p} &\;\equiv\; \mathrm{cap}(p);\;\texttt{\{}~\texttt{let t = read}(p)\texttt{ + 1;}
   ~~\texttt{write}(p, \texttt{t})\texttt{;}~~\texttt{t}~\texttt{\}} \\
\texttt{p--} &\;\equiv\; \mathrm{cap}(p);\;\texttt{\{}~\texttt{let t = read}(p)\texttt{;}
   ~~\texttt{write}(p, \texttt{t - 1})\texttt{;}~~\texttt{t}~\texttt{\}} \\
\texttt{--p} &\;\equiv\; \mathrm{cap}(p);\;\texttt{\{}~\texttt{let t = read}(p)\texttt{ - 1;}
   ~~\texttt{write}(p, \texttt{t})\texttt{;}~~\texttt{t}~\texttt{\}} \\[4pt]
\texttt{p op= e} &\;\equiv\; \mathrm{cap}(p);\;\texttt{\{}~\texttt{let t = read}(p)
   \texttt{ } op \texttt{ } e\texttt{;}~~\texttt{write}(p, \texttt{t})\texttt{;}~~\texttt{t}~\texttt{\}}
\end{align*}

\noindent for $op \in \{+,-,\times,\div,\%\}$ on numbers. The four
increment and decrement forms are expressions: postfix yields the
\emph{old} value, prefix the \emph{new} one. A compound assignment is,
like \ic{=}, a statement form (Section~\ref{sec:assignment}): it cannot
nest inside another expression, and as the last statement of a block it
yields the new value of the place (\ic{\{ x += 2 \}}).

\paragraph{Examples.}
\begin{lstlisting}
a[f()]++;         // { let y_t = f(); let v = a[y_t];
                  //   a[y_t] = v + 1; v }   -- f() runs once
rec.v[g()] += 1;  // g() once; the read and the write use the same element
x.y.z += e;       // x and x.y evaluated once
\end{lstlisting}

\paragraph{Arrays.} On array targets \ic{+=} and \ic{*=} are
\emph{primitive in-place} operations --- they do not desugar through
\ic{+} (which copies) --- and are defined in Section~\ref{sec:arrays}.
On string targets only \ic{+=} is legal and concatenates (in place when
it can, Section~\ref{sec:str-append}). Every other
combination is a type error (Section~\ref{sec:opassign-typing}).

\paragraph{Maps and slices.} A map store \ic{m[k] = v} evaluates \ic{m},
\ic{k} and \ic{v} once each; a slice assignment \ic{a[lo:hi] = b} the same
way (Section~\ref{sec:arrays}).
